\begin{theorem}[Classification in Viaclovsky's coordinates]\hypertarget{tag-032S}{}\label{push:src:E.1.2}
\StatementRecord{push:src:E.1.2}{proved}{032S}{\SeriesRef{push:src:E.1.1}}

Use a local logarithm $c$ of the transported linearization,
$u=e^{-2\pi i c}$, and the same reference section and cusp comparison
on both sides.  Two sphere members with data
$(m,\rho_0,N_0;n,\rho_1,N_1;c)$ and primed data are biholomorphic
if and only if their multiplicities agree and there are
$\alpha\in\{1,-1\}$ and $h=c'-c\in\mathbb Z$ such that
\[
 b'=\alpha b,\qquad
 hN'_0\equiv\rho'_0-\alpha\rho_0\pmod{k},\qquad
 hN'_1\equiv\rho'_1-\alpha\rho_1\pmod{2l}.
\]
Every such biholomorphism covers the identity of $B$.
For fixed discrete data the period identification subgroup is
$2kl\mathbb Z$.  At fixed $u$ all sphere translation data give
exactly $2kl$ classes.  A positive $u$-circuit acts on them by a
single cycle.
\end{theorem}


\begin{proof}
Apply \cref{fam:src:B.3.1} through the integral isomorphism
$P_V$ of \cref{push:src:E.1.1}. Its affine and cusp coordinates are
\[
 r_1=N_0,\quad s_1=-\rho_0,\qquad
 r_2=N_1,\quad s_2=-(\rho_1+N_1),\qquad p=-b.
\]
Thus $p'=\alpha p$ is $b'=\alpha b$, and the first norm congruence is
$hN'_0\equiv\rho'_0-\alpha\rho_0\pmod{k}$.
The second becomes
\[
 hN'_1\equiv\rho'_1-\alpha\rho_1+N'_1-\alpha N_1\pmod{2l}.
\]
The relation $b'=\alpha b$, reduced modulo $4l$, gives
$N'_1\equiv\alpha N_1\pmod{4l}$, since $\gcd(3k,4l)=1$.
It removes the last two terms and yields the stated congruence.
The base-preservation assertion, identification subgroup and number of
classes transfer through this same integral map. It preserves
$u=e^{-2\pi ic}$, so a positive $u$-circuit is still $c\mapsto c-1$
and induces the single cycle proved in \cref{fam:src:B.3.1}.
\end{proof}